Krátké shrnutí: tato sekce nabízí praktický, krok‑za‑krok proces pro zařazení AI do služeb studijní podpory (tutoring, zpětná vazba, diagnostika), rozdílení rolí mezi učitelem a AI‑asistentem a základní etické body, které je třeba před nasazením ošetřit.

Doporučený základní workflow pro nasazení AI v poradenství a podpoře studentů
\begin{enumerate}
  \item Definice cíle a hranic role AI: nejprve stanovte, které konkrétní aktivity má AI provádět (např. generování cvičných materiálů, předběžná kontrola odevzdaných prací, odpovídání na FAQ) a kdy je nutné eskalovat k lidskému pedagogovi. (general background knowledge)
  \item Volba nástrojů a ověření dostupnosti ve školních licencích: preferujte nástroje integrované do univerzitního ekosystému, ověřte dostupnost požadovaných funkcí v školních licencích a možnosti správy přístupů; např. některé funkce rodiny Microsoft (OneNote, Pokrok v matematice, math.microsoft.com) jsou součástí školních licencí a je vhodné předem ověřit jejich dostupnost pro instituci \cite{kb_ai_101_tipu_pdf_04e4a115_page_8_1}. 
  \item Nastavení požadavků na odevzdání a zpracování dat: tam, kde je cílem hodnotit postupy studentů, vyžadujte nahrání postupu (umožňuje to např. nástroj Pokrok v matematice) a nakonfigurujte, jaké výstupy AI smí generovat před pedagogickou revizí \cite{kb_ai_101_tipu_pdf_04e4a115_page_8_1}.
  \item Pilot a ověření end‑to‑end toku: spusťte malý pilot (vybraná skupina kurzů), otestujte generování cvičení, odevzdání, automatické označení problémů a následnou pedagogickou revizi; použijte pilot k doladění pravidel přístupu a retence dat \cite{kb_ai_101_tipu_pdf_04e4a115_page_8_1}.
  \item Integrace diagnostiky do pedagogické práce: využijte předběžné označení chyb AI pro rychlou identifikaci oblastí potřeby remediace a plánování cílených intervencí, přičemž finální hodnocení provádí učitel \cite{kb_ai_101_tipu_pdf_04e4a115_page_8_1}.
  \item Škálování a automatizace rutinních úloh: tam, kde je to vhodné, lze využít nástroje typu Microsoft 365 Copilot v kombinaci s Forms pro rychlé generování cvičných testů nebo kvízů; vždy však ponechte kontrolu nad finálními hodnotami a vysvětlením správných odpovědí \cite{kb_ai_101_tipu_pdf_04e4a115_page_36_1}.
\end{enumerate}

Praktické role a odpovědnosti (kratce)
\begin{itemize}
  \item Učitel / metodik: definice didaktických cílů, ověření validity materiálů generovaných AI, pedagogická revize výsledků a rozhodnutí o závažnosti zásahů do hodnocení. (general background knowledge)
  \item AI‑asistent (systém): automatické generování cvičení, předběžné označení chyb, sumarizace častých nedostatků pro učitele; systém by měl být nastaven tak, aby neprováděl konečné hodnocení bez lidské kontroly \cite{kb_ai_101_tipu_pdf_04e4a115_page_8_1}.
  \item IT / správa dat: zajištění licencí, přístupových práv, zálohování a bezpečného zpracování dat studentů (general background knowledge; viz doporučení ověření funkcí v licencích) \cite{kb_ai_101_tipu_pdf_04e4a115_page_8_1}.
\end{itemize}

Krátký checklist před spuštěním služby (praktické položky)
\begin{itemize}
  \item Ověřit dostupnost požadovaných funkcí (např. nahrávání postupu, převod rukopisu, předopravování) ve školních licencích a nastavit odpovídající přístupy \cite{kb_ai_101_tipu_pdf_04e4a115_page_8_1}.
  \item Nastavit jasná pravidla pro to, kdy AI pouze pomáhá (draft/sugesce) a kdy je vyžadována lidská validace (hodnocení, disciplinární rozhodnutí). (general background knowledge)
  \item Otestovat celý tok (vygenerovat úlohy, simulovat odevzdání, projít AI označeními a pedagogickou revizí) v pilotním běhu \cite{kb_ai_101_tipu_pdf_04e4a115_page_8_1}.
  \item Zajistit záznam o tom, kdo má přístup k výsledkům a jak jsou data uchovávána/mazána (general background knowledge).
\end{itemize}

Základní etické a komunikační požadavky (stručně, doporučené)
\begin{itemize}
  \item Transparentnost vůči studentům: informujte je o tom, kdy a jak AI zasahuje do podpory a hodnocení (general background knowledge).
  \item Ochrana soukromí: minimalizujte sběr citlivých údajů a domluvte retenci dat se správcem IT/fakulty (general background knowledge).
  \item Pedagogická odpovědnost: AI poskytuje pomocnou vrstvu — konečné pedagogické rozhodnutí a oprávnění k hodnocení zůstává u učitele \cite{kb_ai_101_tipu_pdf_04e4a115_page_8_1}.
\end{itemize}

Tip na závěr: pro aktivity zaměřené na seznámení a motivaci studentů (např. krátké workshopy nebo hry přibližující principy AI) lze zvážit využití edukativních prostředí jako Minecraft Education, které se v praxi využívají k demonstraci principů AI a Hodiny kódu; ověřte dostupnost v institucionálních licencích \cite{kb_ai_101_tipu_pdf_04e4a115_page_15_2}.